\subsection{\problemblock{*Summarize} Data Block}\label{sec:block-summarize}

Summary variables are used with surveys to perform tradeoff and sensitivity studies.
A survey systematically changes one or more input variables and calculates a set of
trajectories for each combination, as described in \cref{sec:block-survey}. Summary
variables are calculated after each set of trajectories is complete. When values are
surveyed, the summary variables can therefore be tabulated as functions of the
survey variables.

\Cref{fig:summary-variable-printout} shows summary output for a sounding-rocket
problem run with different payload weights and launch angles. The survey variables
are launch elevation \texttt{qe} and payload weight \texttt{payload}. The summary
variables are burnout velocity \texttt{vel-bo}, maximum altitude or apogee
\texttt{max-alt}, and final range \texttt{range}. Each line represents one trajectory;
ten trajectories were calculated for five payload weights and two launch angles.

\begin{figure}[htbp]
  \centering
  \begin{minipage}{0.78\linewidth}
    \lstinputlisting[style=taosmanual,basicstyle=\ttfamily\scriptsize]{examples/chapter04/summary-printout-example.txt}
  \end{minipage}
  \caption{Summary-variable printout.}
  \label{fig:summary-variable-printout}
\end{figure}

For each trajectory, or each set of trajectories in a multiple-vehicle problem, a
summary variable has one value. The value is calculated from commands in a
\problemblock{*summarize} data block. These commands are math operations of the
same general form used in full table files.

The first summary variable in the figure is defined by

\begin{lstlisting}[style=taosmanual]
*summarize vel-bo
  add vel on segment 3, trajectory 1
\end{lstlisting}

The name \texttt{vel-bo} follows the data-block keyword and is used as the column
heading in the summary printout. Summary-variable names need not be unique.
Summary values are initialized to zero and the operations are then executed. Here the
velocity at the end of segment 3 in trajectory 1 is added to zero. Only one operation
is required. This is typical: summary blocks are usually used to retrieve a specific
trajectory value rather than to perform a long calculation.

The other two summary variables in \cref{fig:summary-variable-printout} are defined
as

\begin{lstlisting}[style=taosmanual]
*summarize max-alt
  add max(alt[1])

*summarize range
  add last(range) trajectory 1
\end{lstlisting}

\taossubhead{Max, Min, First, and Last Special Functions}

The \texttt{max} function searches the complete trajectory for the maximum value of
a variable. The \texttt{last} function retrieves the final value. The corresponding
\texttt{min} and \texttt{first} functions are also available. Because these functions
operate on the entire trajectory, they eliminate the need to provide a segment
number.

TAOS saves trajectory data at the beginning and end of every segment and at every
print interval. When \texttt{max} or \texttt{min} scans the trajectory, it considers
only these saved values. It chooses the largest or smallest saved value and uses that
sample and its two neighbors to fit a parabola and estimate the actual extremum. The
accuracy therefore depends on the print interval. If the variable changes rapidly, the
print interval may need to be reduced.

When a special function is not used, segment and trajectory numbers are normally
required. Segment numbers follow the words \texttt{on segment}; trajectory numbers
follow \texttt{trajectory}. A trajectory number can alternatively be written as a
subscript in square brackets, as in \texttt{max(alt[1])}.

\taossubhead{Math Operations}

All preceding examples use \texttt{add}, but the complete set of operations is given
in \cref{tab:summary-math-operations}. The operations from \texttt{add} through
\texttt{iexp} require an additional value and operate on both that value and the
current summary. The additional value can be a trajectory output variable or a
constant. The operations from \texttt{abs} through \texttt{atan} act only on the
current summary and require no additional value.

\begin{longtable}{@{}>{\ttfamily}lp{0.29\linewidth}>{\ttfamily}lp{0.32\linewidth}@{}}
  \caption{Summary-variable math operations.}
  \label{tab:summary-math-operations}\\
  \toprule
  \normalfont Operation & Description & \normalfont Operation & Description \\
  \midrule
  \endfirsthead
  \toprule
  \normalfont Operation & Description & \normalfont Operation & Description \\
  \midrule
  \endhead
  add  & Add a value to the summary value & sqrt & Square root of summary value \\
  sub  & Subtract a value from the summary value & ln & Natural logarithm of summary value \\
  mult & Multiply a value by the summary value & log & Base-10 logarithm of summary value \\
  div  & Divide the summary value by a supplied value & e & $e$ raised to the summary value \\
  idiv & Divide a supplied value by the summary value & sin & Sine of summary value in degrees \\
  exp  & Raise the summary value to a supplied power & cos & Cosine of summary value in degrees \\
  iexp & Raise a supplied value to the summary-value power & tan & Tangent of summary value in degrees \\
  abs  & Absolute value of summary value & asin & Inverse sine of summary value in degrees \\
  neg  & Negate summary value & acos & Inverse cosine of summary value in degrees \\
  sqr  & Square summary value & atan & Inverse tangent of summary value in degrees \\
  \bottomrule
\end{longtable}

A block may contain any number of operations, although most summary variables need
only one or two. The following three-operation example calculates the difference in
final range between two trajectories and converts the result from nautical miles to
kilometers:

\begin{lstlisting}[style=taosmanual]
*summarize delta-rng
  add last(range[1])
  sub last(range[2])
  mult 1.852
\end{lstlisting}

The summary starts at zero. The first operation adds the final range of trajectory 1,
the second subtracts the final range of trajectory 2, and the third converts units.

The preceding example uses \texttt{last}, so no segment numbers are needed. The
next example retrieves values from within two trajectories:

\begin{lstlisting}[style=taosmanual]
*summarize delta-alt
  add alt on segment 4, trajectory 1
  sub alt on segment 11, trajectory 2
  div 3280.84
\end{lstlisting}

It calculates, in kilometers, the difference between altitude at the end of segment 4
in trajectory 1 and altitude at the end of segment 11 in trajectory 2. Segment numbers
are required because none of the special functions is used. The trajectory numbers
are entered with the \texttt{trajectory} keyword to illustrate that form.
